\chapter{Sums of Random Variables and Convolutions}

The previous chapter developed tools for general transformations of random
variables. We now specialize to addition: $S = X + Y$. Given the individual
distributions of independent random variables $X$ and $Y$, what is the
distribution of their sum $S$?

The answer is the \textbf{convolution} of their distributions. We derive the
convolution formula for discrete and continuous variables, examine families
closed under addition (such as Poisson, binomial, and normal), and see how sums
of independent variables behave.

\chapterline

\section{The Question: What Is the Distribution of a Sum?}
\label{ch17:question}

Start with the simplest possible story.

\begin{intuitionbox}
You roll two fair four-sided dice (a $\Unif\set{1,2,3,4}$ each). You do
not care about the individual faces; you care about the \emph{total}
$S = X + Y$. What is $\P(S = 5)$?

You already know how to do this by brute force: list the $16$ equally
likely outcomes $(x,y)$ and count those with $x+y=5$. They are
$(1,4),(2,3),(3,2),(4,1)$ -- four of them -- so $\P(S=5)=4/16=1/4$.

Stare at that list. Each pair has the form $(k,\,5-k)$. We summed over all
the ways the two pieces can ``conspire'' to produce the target total $5$:
the first die is $k$, the second die is $5-k$, and we add up the
probability of each such split. That summing-over-all-splits is exactly
what a convolution is. The rest of this chapter is just that idea, made
precise and pushed into the continuous world.
\end{intuitionbox}

Before naming the tool, we should be clear about \emph{when} it applies and
what could go wrong.

\begin{focusbox}
The distribution of $S = X+Y$ depends, in general, on the \emph{joint}
distribution of $(X,Y)$, not just on the two marginals. Two random
variables with the same individual distributions can produce very
different sums depending on how they are coupled.

\textbf{Extreme example.} Let $X \sim \Unif\set{0,1}$ (a fair coin coded
$0/1$).
\begin{itemize}
  \item If $Y = X$ (perfectly dependent), then $S = 2X \in \set{0,2}$,
        each with probability $1/2$. It can never equal $1$.
  \item If $Y = 1-X$ (perfectly anti-dependent), then $S \equiv 1$ always.
  \item If $Y$ is an independent copy of $X$, then
        $S \in \set{0,1,2}$ with probabilities $\tfrac14,\tfrac12,\tfrac14$.
\end{itemize}
Same marginals, three completely different sums. So the clean formulas of
this chapter will require an assumption, and that assumption is
\textbf{independence}.
\end{focusbox}

Throughout, unless we say otherwise, $X \indep Y$.

\begin{recoverybox}
\textbf{Pause and think.} Suppose $X$ and $Y$ are independent fair coin flips, coded as $0$ (tails) or $1$ (heads). What are the possible values of $S = X + Y$, and how many ways can you land on $S = 1$?

\emph{Check:} The possible totals are $0, 1, 2$. To land on $S = 1$, the pieces must split as $(X=0, Y=1)$ or $(X=1, Y=0)$. Each pair has probability $\tfrac14$, so $\P(S=1) = \tfrac14 + \tfrac14 = \tfrac12$. You just computed a discrete convolution in your head: you held the target sum $s=1$ fixed and summed the probabilities of all valid splits $(k,\,1-k)$.
\end{recoverybox}

\section{Discrete Convolution: Summing Over All the Ways}
\label{ch17:discrete-conv}

Suppose $X$ and $Y$ are independent integer-valued random variables with
probability mass functions (PMFs) $p_X$ and $p_Y$ -- rules assigning each
possible discrete integer its exact probability. We want the PMF of
$S = X+Y$. Fix a target value $s$. The event $\set{S=s}$ is the union,
over every integer $k$, of the event ``$X=k$ \emph{and} $Y=s-k$.'' These
are disjoint (different $k$ means a different value of $X$), so we add
their probabilities, and independence lets us factor each one:
\[
\P(S=s) \;=\; \sum_{k} \P(X=k,\,Y=s-k)
            \;=\; \sum_{k} \P(X=k)\,\P(Y=s-k).
\]
That is the whole derivation. We give it a name.

\begin{definition}[Discrete convolution]
\label{ch17:def-discrete-conv}
For independent integer-valued $X,Y$ with PMFs $p_X,p_Y$, the PMF of
$S=X+Y$ is the \emph{convolution} $p_X * p_Y$:
\[
p_S(s) \;=\; (p_X * p_Y)(s) \;=\; \sum_{k=-\infty}^{\infty} p_X(k)\,p_Y(s-k).
\]
\end{definition}

\begin{happeningbox}
Read the formula as a recipe. To land the total at $s$, you must split $s$
into a part contributed by $X$ (call it $k$) and the leftover $s-k$
contributed by $Y$. There are many possible splits; you weight each split
by how likely it is ($p_X(k)\,p_Y(s-k)$) and add up. The index $k$ slides
across all values while $s-k$ slides in the opposite direction -- one
counter goes up exactly as the other goes down so their sum is pinned at
$s$. That ``one up, one down'' sliding is the visual signature of
convolution, and it is why engineers describe convolution as ``flip one
sequence and slide it past the other.''
\end{happeningbox}

\begin{connectionbox}
If you have met convolution in the companion calculus book (or in signal
processing), this is the \emph{same} operation applied to PMFs.
In particular, multiplying two polynomials is precisely a convolution of their
coefficients: if
$A(z)=\sum_k a_k z^k$ and $B(z)=\sum_k b_k z^k$, then the coefficient of
$z^s$ in the product $A(z)B(z)$ is exactly $\sum_k a_k b_{s-k}$. We will exploit this
through the idea of \emph{generating functions}: turn a PMF into a
polynomial (or power series), and adding independent random variables
becomes \emph{multiplying} their generating functions. That turns hard
convolutions into easy products -- the central idea of the next chapter.
\end{connectionbox}

\subsection{Worked example: two four-sided dice}

Let $p_X(k)=p_Y(k)=\tfrac14$ for $k\in\set{1,2,3,4}$ and $0$ otherwise.
By Definition~\ref{ch17:def-discrete-conv},
\[
p_S(s) \;=\; \sum_{k=1}^{4} \tfrac14 \cdot p_Y(s-k)
        \;=\; \tfrac{1}{16}\,\#\set{k\in\set{1,2,3,4} : s-k \in \set{1,2,3,4}}.
\]
So $p_S(s)$ is just $1/16$ times the number of valid splits. Counting them
for $s=2,\dots,8$ gives $1,2,3,4,3,2,1$ -- the famous triangular
``pyramid'' of dice sums. Figure~\ref{ch17:fig-dice} plots it.

\begin{figure}[H]
\centering
\begin{tikzpicture}
\begin{axis}[
  width=12cm, height=6cm,
  ybar, bar width=14pt,
  xlabel={total $s$}, ylabel={$\P(S=s)$},
  xtick={2,3,4,5,6,7,8},
  ymin=0, ymax=0.30,
  enlarge x limits=0.12,
  nodes near coords, every node near coord/.append style={font=\scriptsize},
  point meta=explicit symbolic,
]
\addplot[fill=distblue!60, draw=distblue] coordinates {
  (2,0.0625)[$\tfrac{1}{16}$] (3,0.125)[$\tfrac{2}{16}$]
  (4,0.1875)[$\tfrac{3}{16}$] (5,0.25)[$\tfrac{4}{16}$]
  (6,0.1875)[$\tfrac{3}{16}$] (7,0.125)[$\tfrac{2}{16}$]
  (8,0.0625)[$\tfrac{1}{16}$]
};
\end{axis}
\end{tikzpicture}
\caption{The PMF of the sum of two fair four-sided dice. Two flat
(uniform) distributions convolve into a triangle. This smoothing --
flat becoming peaked -- illustrates the mechanism underlying the central limit theorem.}
\label{ch17:fig-dice}
\end{figure}

\begin{aestheticbox}
Notice what convolution did to the \emph{shape}. Each die alone is flat:
every face equally likely. Their sum is a triangle, peaked in the middle.
Convolve a third die in and you get a bell-ish bump; a fourth, smoother
still. Adding independent copies \emph{pulls probability toward the center}
and rounds off the corners. Keep this picture in your pocket: it is the
mechanism behind the bell curve, and we will cash it out fully when we
reach the central limit theorem.
\end{aestheticbox}

\subsection{A nontrivial discrete example: two Poissons}

Convolution is not just for finite supports. Let
$X \sim \Pois(\lambda)$ and $Y\sim\Pois(\mu)$ be independent, supported on
$\set{0,1,2,\dots}$. For $s\ge 0$,
\[
\P(S=s) = \sum_{k=0}^{s} e^{-\lambda}\frac{\lambda^k}{k!}\,
          e^{-\mu}\frac{\mu^{\,s-k}}{(s-k)!}
        = e^{-(\lambda+\mu)} \sum_{k=0}^{s} \frac{\lambda^k \mu^{\,s-k}}{k!\,(s-k)!}.
\]
Multiply and divide by $s!$ inside the sum to manufacture binomial
coefficients:
\[
\P(S=s) = \frac{e^{-(\lambda+\mu)}}{s!}
          \sum_{k=0}^{s} \binom{s}{k}\lambda^k \mu^{\,s-k}
        = e^{-(\lambda+\mu)}\,\frac{(\lambda+\mu)^s}{s!},
\]
where the last step is the binomial theorem. The right-hand side is the
PMF of $\Pois(\lambda+\mu)$. We have just proved one of the most-used
facts in applied probability.

\begin{theorem}[Poissons add]
\label{ch17:thm-pois-add}
If $X\sim\Pois(\lambda)$ and $Y\sim\Pois(\mu)$ are independent, then
$X+Y \sim \Pois(\lambda+\mu)$.
\end{theorem}

\begin{intuitionbox}
This should feel \emph{right} before you even see the algebra. The Poisson
counts rare events over an interval at rate $\lambda$. If one phone line
gets calls at rate $\lambda$ and another, independently, at rate $\mu$,
then the combined switchboard gets calls at rate $\lambda+\mu$. Merging
two independent streams of rare events just adds their rates. The
convolution computation is merely the bookkeeping that confirms the story.
\end{intuitionbox}

\section{Continuous Convolution: From Sums to Integrals}
\label{ch17:cont-conv}

Now let $X$ and $Y$ be independent \emph{continuous} random variables with
probability density functions (PDFs or densities) $f_X$ and $f_Y$ --
curves where probability corresponds to area under the curve rather than
individual point heights. We want the density of $S = X+Y$. The route
that actually works runs through the cumulative distribution function
(CDF, the running probability $F_S(s) = \P(S \le s)$), because the
probability of any single point is zero and cannot be added the way PMFs
were.

\subsection{Derivation via the CDF}

Fix a target number $s$. We compute $F_S(s) = \P(X+Y \le s)$ by integrating
the joint density over the half-plane $\set{(x,y): x+y\le s}$. By
independence the joint density factors, $f_{X,Y}(x,y)=f_X(x)f_Y(y)$:
\[
F_S(s) = \iint_{x+y\le s} f_X(x)\,f_Y(y)\,\dd x\,\dd y
       = \int_{-\infty}^{\infty} f_X(x)
         \left( \int_{-\infty}^{s-x} f_Y(y)\,\dd y \right) \dd x.
\]
The inner integral is $F_Y(s-x)$. Now differentiate in $s$ to recover the
density. Differentiation passes under the outer integral, and
$\frac{\dd}{\dd s}F_Y(s-x)=f_Y(s-x)$:
\[
f_S(s) = \frac{\dd}{\dd s}F_S(s)
       = \int_{-\infty}^{\infty} f_X(x)\,f_Y(s-x)\,\dd x.
\]

\begin{figure}[H]
\centering
\begin{tikzpicture}[scale=1.05]
  \fill[sourceblue!15] (-0.6,-0.6) -- (-0.6,3.4) -- (3.4,-0.6) -- cycle;
  \fill[quotientorange!35] (-0.6,3.4) -- (-0.6,3.8) -- (3.8,-0.6) -- (3.4,-0.6) -- cycle;

  \draw[->, thick, gray!70] (-1.0,0) -- (4.6,0) node[right, font=\small] {$x$};
  \draw[->, thick, gray!70] (0,-1.0) -- (0,4.6) node[above, font=\small] {$y$};

  \draw[sourceblue!80!black, very thick] (-0.8,3.6) -- (3.6,-0.8)
    node[pos=0.10, above right=1pt, font=\small\bfseries, sourceblue!90!black] {$x + y = s$};

  \draw[quotientorange!90!black, thick, dashed] (-0.8,4.0) -- (4.0,-0.8);

  \fill[sourceblue!80!black] (2.8,0) circle (2pt) node[below=3pt, font=\footnotesize] {$(s, 0)$};
  \fill[sourceblue!80!black] (0,2.8) circle (2pt) node[left=3pt, font=\footnotesize] {$(0, s)$};

  \fill[densityred] (1.2, 1.6) circle (2.2pt);
  \draw[densely dotted, gray!60] (1.2,0) -- (1.2,1.6) -- (0,1.6);
  \node[above right=2pt, font=\footnotesize\bfseries, densityred!90!black] at (1.2,1.6) {$(x,\, s - x)$};
  \node[below, font=\scriptsize, gray!80!black] at (1.2,0) {$x$};
  \node[left, font=\scriptsize, gray!80!black] at (0,1.6) {$y=s-x$};

  \draw[->, >=stealth, very thick, quotientorange!90!black] (2.1, 2.1) -- (2.8, 2.8)
    node[above right=1pt, font=\footnotesize\bfseries, align=left] {sweep as\\$s$ grows};

  \node[sourceblue!90!black, font=\small\bfseries] at (0.8, 0.75) {$\set{X + Y \le s}$};
  \node[sourceblue!80!black, font=\scriptsize] at (0.8, 0.35) {CDF area $F_S(s)$};

  \node[quotientorange!90!black, font=\scriptsize\bfseries, align=left] at (3.8, 1.0)
    {slice of width $\dd s$:\\density $f_S(s)\,\dd s$};
  \draw[->, >=stealth, thin, quotientorange!90!black] (3.1, 1.0) -- (2.3, 0.8);
\end{tikzpicture}
\caption{The geometry of sums in the $(x,y)$ plane. To find the cumulative probability $F_S(s) = \P(X+Y \le s)$, we integrate the joint density $f_{X,Y}(x,y)$ over the half-plane below the diagonal line $x+y=s$ (shaded blue). Differentiating with respect to $s$ isolates the probability in a thin diagonal strip of width $\dd s$ along the line $y = s-x$ (shaded orange). The continuous convolution formula $\int f_X(x) f_Y(s-x)\,\dd x$ integrates the joint density along this diagonal slice.}
\label{ch17:fig-plane}
\end{figure}

\begin{definition}[Continuous convolution]
\label{ch17:def-cont-conv}
For independent continuous $X,Y$ with densities $f_X,f_Y$, the density of
$S=X+Y$ is the \emph{convolution}
\[
f_S(s) = (f_X * f_Y)(s) = \int_{-\infty}^{\infty} f_X(x)\,f_Y(s-x)\,\dd x.
\]
\end{definition}

\begin{connectionbox}
\textbf{The bivariate transformation view from Chapter~16.}
In Chapter~16, we saw how to find the distribution of any smooth bivariate
transformation $(X, Y) \mapsto (U, V)$ using the Jacobian determinant. We
can derive continuous convolution the exact same way.
Define the transformed coordinates:
\[
S = X + Y, \qquad U = X,
\]
where $U$ is a convenient auxiliary coordinate that preserves the
two-dimensional area. Solving backwards gives the inverse map $X = U$ and
$Y = S - U$. The Jacobian matrix of this transformation is
\[
J = \begin{pmatrix} \frac{\partial x}{\partial s} & \frac{\partial x}{\partial u} \\ \frac{\partial y}{\partial s} & \frac{\partial y}{\partial u} \end{pmatrix}
  = \begin{pmatrix} 0 & 1 \\ 1 & -1 \end{pmatrix},
\]
which has determinant $\det(J) = 0\cdot(-1) - 1\cdot 1 = -1$, so
$\abs{\det(J)} = 1$. By the bivariate change-of-variables formula from
Chapter~16, the joint density of $(S, U)$ is
\[
f_{S,U}(s,u) = f_{X,Y}(u, s-u)\,\abs{\det(J)} = f_X(u)\,f_Y(s-u),
\]
where independence allowed us to factor the joint density. To recover the
marginal density of $S$ alone, we integrate out the auxiliary variable $u$:
\[
f_S(s) = \int_{-\infty}^{\infty} f_{S,U}(s,u)\,\dd u
       = \int_{-\infty}^{\infty} f_X(u)\,f_Y(s-u)\,\dd u.
\]
Renaming the dummy variable $u$ to $x$ reproduces
Definition~\ref{ch17:def-cont-conv} immediately. The CDF method and the
Jacobian method are two perspectives on the exact same geometry.
\end{connectionbox}

\begin{connectionbox}
Compare Definitions~\ref{ch17:def-discrete-conv} and
\ref{ch17:def-cont-conv}: the sum $\sum_k$ has become an integral
$\int \dd x$, and the PMFs have become densities. This is exactly the
discrete-to-continuous translation you saw for expectation in the calculus
companion -- a Riemann sum passing to its integral.
\end{connectionbox}

\begin{whybox}
Why does differentiating ``under the integral sign'' give the density?
Because for nice (continuous, integrable) $f_X,f_Y$ the map
$s\mapsto \int f_X(x)F_Y(s-x)\dd x$ is differentiable and its derivative is
$\int f_X(x)\,f_Y(s-x)\dd x$, by the Leibniz rule for differentiating an
integral with respect to a parameter. The regularity conditions hold for
every density we use in this book; the rigorous statement lives in the
measure-theory chapters. For now, trust that the CDF is smooth enough to
differentiate term by term.
\end{whybox}

\subsection{Worked example: sum of two uniforms is a triangle}

Let $X,Y$ be independent $\Unif(0,1)$, so $f_X=f_Y=\ind_{[0,1]}$. Then
\[
f_S(s) = \int_{-\infty}^{\infty} \ind_{[0,1]}(x)\,\ind_{[0,1]}(s-x)\,\dd x
       = \int_0^1 \ind_{[0,1]}(s-x)\,\dd x.
\]
The integrand is $1$ exactly when $0\le s-x\le 1$, i.e. when
$s-1\le x\le s$. Intersecting that with $[0,1]$:
\begin{itemize}
  \item For $0\le s\le 1$: the overlap is $[0,s]$, length $s$. So $f_S(s)=s$.
  \item For $1\le s\le 2$: the overlap is $[s-1,1]$, length $2-s$. So $f_S(s)=2-s$.
  \item Otherwise $f_S(s)=0$.
\end{itemize}
The result is the \emph{triangular} (or ``tent'') density on $[0,2]$,
peaking at $s=1$. Figure~\ref{ch17:fig-tent} illustrates the entire flip, slide, and accumulation process.

\begin{figure}[H]
\centering
\begin{tikzpicture}[scale=0.92]
  % Lane 1: Sliding Snapshots at s = 0.5, s = 1.0, s = 1.5
  \begin{scope}[shift={(0,4.2)}]
    \draw[draw=gray!30, fill=gray!3, rounded corners=4pt] (-1.0,-0.8) rectangle (13.8,2.6);
    \node[anchor=north, font=\small\bfseries, structurepurple!90!black] at (6.4,2.5) {1. The Flip-and-Slide Mechanism: $f_Y(s-x)$ sliding over $f_X(x)$};

    % Snapshot A: s = 0.5
    \begin{scope}[shift={(0,0)}]
      \draw[->, thick, gray!70] (-1.2,0) -- (3.2,0) node[right, font=\tiny] {$x$};
      \draw[->, thick, gray!70] (0,0) -- (0,1.5);
      \draw[sourceblue!80!black, thick] (0,0) -- (0,1) -- (2,1) -- (2,0);
      \draw[densityred!80!black, thick, dashed] (-1,0) -- (-1,1) -- (1,1) -- (1,0);
      \fill[quotientorange!40, draw=quotientorange!90!black, thick] (0,0) rectangle (1,1);
      \node[font=\scriptsize\bfseries] at (1,-0.35) {$s = 0.5$};
      \node[font=\tiny\bfseries, quotientorange!90!black, align=center] at (0.5,0.5) {area\\$0.5$};
    \end{scope}

    % Snapshot B: s = 1.0
    \begin{scope}[shift={(4.8,0)}]
      \draw[->, thick, gray!70] (-0.5,0) -- (3.5,0) node[right, font=\tiny] {$x$};
      \draw[->, thick, gray!70] (0,0) -- (0,1.5);
      \draw[sourceblue!80!black, thick] (0,0) -- (0,1) -- (2,1) -- (2,0);
      \draw[densityred!80!black, thick, dashed] (0,0) -- (0,1.05) -- (2,1.05) -- (2,0);
      \fill[quotientorange!50, draw=quotientorange!90!black, thick] (0,0) rectangle (2,1);
      \node[font=\scriptsize\bfseries] at (1,-0.35) {$s = 1.0$};
      \node[font=\tiny\bfseries, quotientorange!90!black, align=center] at (1,0.5) {peak area\\$1.0$};
    \end{scope}

    % Snapshot C: s = 1.5
    \begin{scope}[shift={(9.6,0)}]
      \draw[->, thick, gray!70] (-0.5,0) -- (3.5,0) node[right, font=\tiny] {$x$};
      \draw[->, thick, gray!70] (0,0) -- (0,1.5);
      \draw[sourceblue!80!black, thick] (0,0) -- (0,1) -- (2,1) -- (2,0);
      \draw[densityred!80!black, thick, dashed] (1,0) -- (1,1) -- (3,1) -- (3,0);
      \fill[quotientorange!40, draw=quotientorange!90!black, thick] (1,0) rectangle (2,1);
      \node[font=\scriptsize\bfseries] at (1,-0.35) {$s = 1.5$};
      \node[font=\tiny\bfseries, quotientorange!90!black, align=center] at (1.5,0.5) {area\\$0.5$};
    \end{scope}
  \end{scope}

  % Transmission arrows down to Lane 2
  \draw[->, >=stealth, thick, quotientorange!90!black] (1.0, 3.3) -- (3.75, 2.1)
    node[midway, below left, font=\scriptsize\bfseries] {overlap $\frac{1}{2}$};
  \draw[->, >=stealth, thick, quotientorange!90!black] (5.8, 3.3) -- (5.5, 2.1)
    node[midway, right=3pt, font=\scriptsize\bfseries] {overlap $1$};
  \draw[->, >=stealth, thick, quotientorange!90!black] (10.6, 3.3) -- (7.25, 2.1)
    node[midway, below right, font=\scriptsize\bfseries] {overlap $\frac{1}{2}$};

  % Lane 2: Accumulated Output Triangle Density
  \begin{scope}[shift={(2.0,-1.4)}]
    \draw[draw=gray!30, fill=gray!3, rounded corners=4pt] (-1.0,-0.8) rectangle (9.2,3.4);
    \node[anchor=north, font=\small\bfseries, densityred!90!black] at (4.1,3.3) {2. Output Sum Density: Triangular Distribution $f_S(s) = (f_X * f_Y)(s)$};

    \draw[->, thick, gray!70] (-0.4,0) -- (8.2,0) node[right, font=\scriptsize] {$s$};
    \draw[->, thick, gray!70] (0,-0.2) -- (0,2.35) node[above, font=\scriptsize] {$f_S(s)$};

    \fill[densityred!15] (0,0) -- (3.5,1.8) -- (7.0,0) -- cycle;
    \draw[densityred, very thick] (-0.4,0) -- (0,0) -- (3.5,1.8) -- (7.0,0) -- (7.8,0);

    \foreach \s/\pos in {0/0, 0.5/1.75, 1.0/3.5, 1.5/5.25, 2.0/7.0} {
      \draw[gray!70] (\pos, 0.08) -- (\pos, -0.08) node[below, font=\scriptsize] {$\s$};
    }
    \draw[gray!70] (0.08, 1.8) -- (-0.08, 1.8) node[left, font=\scriptsize] {$1.0$};
    \draw[gray!70] (0.08, 0.9) -- (-0.08, 0.9) node[left, font=\scriptsize] {$0.5$};

    \draw[dashed, quotientorange!80!black] (1.75,0) -- (1.75,0.9) -- (0,0.9);
    \fill[quotientorange!90!black] (1.75,0.9) circle (2.2pt) node[above left=-1pt, font=\tiny\bfseries] {$(0.5, 0.5)$};

    \draw[dashed, quotientorange!80!black] (3.5,0) -- (3.5,1.8) -- (0,1.8);
    \fill[densityred] (3.5,1.8) circle (2.4pt) node[above, font=\tiny\bfseries] {Peak $(1.0, 1.0)$};

    \draw[dashed, quotientorange!80!black] (5.25,0) -- (5.25,0.9);
    \fill[quotientorange!90!black] (5.25,0.9) circle (2.2pt) node[above right=-1pt, font=\tiny\bfseries] {$(1.5, 0.5)$};
  \end{scope}
\end{tikzpicture}
\caption{The sliding-window geometry of convolution. (Top) The flipped kernel $f_Y(s-x)$ slides across $f_X(x)$ as $s$ advances. At each shift $s$, the shaded overlap represents the product integral $\int f_X(x) f_Y(s-x)\,\dd x$. (Bottom) Accumulating this overlap area across all shifts $s \in [0,2]$ yields the continuous triangular density $f_S(s)$, turning sharp steps into a smooth peaked tent.}
\label{ch17:fig-tent}
\end{figure}

\begin{walkbox}
\textbf{Walking question.} Before reading on: what do you expect the
density of the sum of \emph{three} independent $\Unif(0,1)$ variables to
look like? Sketch it.

\emph{Answer.} Convolving the tent with one more box rounds the two
straight sides into curves; you get a smooth bell-shaped piece made of
three parabolic arcs (the Irwin--Hall density for $n=3$), supported on
$[0,3]$ and peaking at $1.5$. Each extra uniform you add makes the result
smoother and more bell-like. By $n=12$ the sum is so close to a normal
that the old ``add twelve uniforms and subtract six'' trick was a standard
way to fake a standard normal before fast Gaussian samplers existed.
\end{walkbox}

\begin{historybox}
The word ``convolution'' (German \emph{Faltung}, literally ``a folding'')
spread through analysis and engineering in the early twentieth century to
describe reflecting one signal and sliding it across another, exactly like
the geometric sliding in Figure~\ref{ch17:fig-tent}. But the mathematical
challenge of adding independent variables began two centuries earlier.

In the eighteenth and nineteenth centuries, astronomers and surveyors
combining repeated telescope sightings of planets and stars faced an urgent
version of this problem. Every observation carried unavoidable random errors,
and calculating the most probable orbit required the probability distribution
of the total accumulated error -- a sum of many independent random variables.
Armed only with classical calculus, Taylor series, and generating polynomials
(which Abraham de Moivre had used in 1738 to count dice splits by tracking
powers of a dummy variable $z$), and lacking modern measure theory, functional
analysis, or digital computers, mathematicians found that evaluating
multi-dimensional convolution integrals directly hit an intractable
computational wall: each additional measurement added another nested integral.

In 1810, Pierre-Simon Laplace made an audacious leap: he
bypassed the brutal nested convolution integrals by converting the error
densities into imaginary exponentials -- what we now call
\emph{characteristic functions} (Fourier transforms that wrap probability
densities around the complex unit circle). In this transformed space,
convolutions turn into ordinary function multiplication.
However, Laplace manipulated complex integrals with bold, heuristic steps,
shifting integration contours through the complex plane and exchanging
limits without justification. Rigorous mathematicians, led by
Augustin-Louis Cauchy, Sim\'eon Denis Poisson (in the 1830s), and later
Pafnuty Chebyshev, were deeply skeptical,
demanding proof that these transform shortcuts and normal limits were
mathematically sound rather than lucky approximations. It took more than a
century -- culminating in the work of Aleksandr Lyapunov (1901) and Paul
L\'evy (1925) -- to establish the rigorous analytic foundations of
characteristic functions.

Once placed on firm ground, convolutions and their transform shortcuts
unlocked the Central Limit Theorem, the theory of diffusion and Brownian
motion, error analysis across all experimental sciences, and the entire
field of digital signal processing. Even today, new questions continue to
emerge. In classical probability, numbers commute:
$X+Y = Y+X$. In the 1980s, Dan-Virgil Voiculescu asked what happens when we
add random variables that do \emph{not} commute, such as large random
matrices arising in quantum mechanics and deep neural networks. This
created \emph{free probability theory}, where classical convolution is
replaced by ``free convolution.'' Today, finding optimal convergence rates
for multi-dimensional convolutions (the multivariate Berry--Esseen theorem)
and analyzing convolutions of dependent or heavy-tailed variables on massive
graphs remain active frontiers of mathematical research.
\end{historybox}

\section{Sums Within a Family: Closure Under Addition}
\label{ch17:closure}

Some families of distributions have a remarkable stability property: add
two members and you stay inside the family, with parameters that simply
combine. We already proved it for the Poisson. Here is the roster you
should know cold.

\subsection{Binomial plus binomial (same \texorpdfstring{$p$}{p})}

\begin{theorem}[Binomials add when $p$ matches]
\label{ch17:thm-bin-add}
If $X\sim\Bin(m,p)$ and $Y\sim\Bin(n,p)$ are independent \emph{with the
same success probability $p$}, then $X+Y\sim\Bin(m+n,p)$.
\end{theorem}

\begin{whybox}
The slick proof needs no algebra. A $\Bin(m,p)$ counts successes in $m$
independent coin flips with bias $p$; a $\Bin(n,p)$ counts successes in
another $n$ such flips. Lay all $m+n$ flips in a row -- they are
independent and all have bias $p$ -- and $X+Y$ is the total number of
successes among $m+n$ flips, which is by definition $\Bin(m+n,p)$. The
constraint ``same $p$'' is essential: if the coins had different biases,
the pooled flips would not be identically distributed and the count would
not be binomial.
\end{whybox}

If you prefer to see the convolution machinery confirm it, the algebraic
heart is Vandermonde's identity
$\sum_{k}\binom{m}{k}\binom{n}{s-k}=\binom{m+n}{s}$, which collects the
ways to get $s$ total successes from the two groups.

\subsection{Normal plus normal}

\begin{theorem}[Normals add]
\label{ch17:thm-normal-add}
If $X\sim\Normal(\mu_1,\sigma_1^2)$ and $Y\sim\Normal(\mu_2,\sigma_2^2)$
are independent, then
\[
X+Y \sim \Normal\!\left(\mu_1+\mu_2,\;\sigma_1^2+\sigma_2^2\right).
\]
The means add and the \emph{variances} add (so the standard deviations add
in quadrature, $\sqrt{\sigma_1^2+\sigma_2^2}$, not linearly -- that is,
standard deviations combine like orthogonal legs of a right triangle
under the Pythagorean theorem).
\end{theorem}

\begin{recoverybox}
\textbf{Pause and think.} Suppose $X \sim \Normal(0, 9)$ and $Y \sim \Normal(0, 16)$ are independent. What is the standard deviation of their sum $X+Y$?

\emph{Check:} Variances add: $\Var(X+Y) = 9 + 16 = 25$. Taking the square root gives $\SD(X+Y) = \sqrt{25} = 5$. Notice this is the classic $3\text{--}4\text{--}5$ right triangle! The individual standard deviations are $3$ and $4$, and because the variables are independent, their combined spread is $\sqrt{3^2 + 4^2} = 5$, not $3 + 4 = 7$.
\end{recoverybox}

\begin{whybox}
Here is the convolution computation for the clean case
$\mu_1=\mu_2=0$ (a shift handles general means). With
$f_X(x)=\frac{1}{\sqrt{2\pi}\sigma_1}e^{-x^2/2\sigma_1^2}$ and similarly
for $Y$,
\[
f_S(s) = \frac{1}{2\pi\sigma_1\sigma_2}\int_{-\infty}^{\infty}
   \exp\!\left(-\frac{x^2}{2\sigma_1^2}-\frac{(s-x)^2}{2\sigma_2^2}\right)\dd x.
\]
The exponent is a quadratic in $x$. ``Completing the square'' in $x$
splits it into a part that depends only on $s$ and a part that integrates
to a constant Gaussian integral over $x$. Carrying out the algebra (group
the $x^2$ terms, the $x$ terms, the constant terms; the $x$-integral
becomes a standard $\int e^{-a(x-b)^2}\dd x=\sqrt{\pi/a}$) leaves
\[
f_S(s) = \frac{1}{\sqrt{2\pi(\sigma_1^2+\sigma_2^2)}}\,
         \exp\!\left(-\frac{s^2}{2(\sigma_1^2+\sigma_2^2)}\right),
\]
which is exactly the $\Normal(0,\sigma_1^2+\sigma_2^2)$ density. This follows
because a Gaussian times a Gaussian, integrated, yields another Gaussian:
the product of two exponentials of quadratics is itself the exponential of a
quadratic, and completing the square isolates the integration variable. (Using
characteristic functions in the next chapter, the same fact takes one
line.)
\end{whybox}

\begin{aestheticbox}
The normal family is the \emph{fixed point} of convolution. Almost every
other family changes shape when you add copies (uniform becomes triangle
becomes bell), forever crawling toward the normal. The normal is already
there. It does not move. That self-reproducing stability is the deep
reason the normal is the universal limit: it is the only ``shape'' that
adding cannot disturb (among distributions with finite variance). We make
this precise in the central limit theorem.
\end{aestheticbox}

\subsection{Gamma plus gamma (same rate)}

The Gamma family with shape $\alpha$ and rate $\beta$ has density
$f(x)=\frac{\beta^\alpha}{\Gamma(\alpha)}x^{\alpha-1}e^{-\beta x}$ for
$x>0$. It too is closed under addition when the rate matches.

\begin{theorem}[Gammas add when the rate matches]
\label{ch17:thm-gamma-add}
If $X\sim\Gam(\alpha_1,\beta)$ and $Y\sim\Gam(\alpha_2,\beta)$ are
independent with the same rate $\beta$, then
$X+Y\sim\Gam(\alpha_1+\alpha_2,\beta)$.
\end{theorem}

\begin{whybox}
Convolve the densities for $s>0$ (constants pulled out, $\beta$ shared):
\[
f_S(s) = \frac{\beta^{\alpha_1+\alpha_2}}{\Gamma(\alpha_1)\Gamma(\alpha_2)}
         e^{-\beta s}\int_0^s x^{\alpha_1-1}(s-x)^{\alpha_2-1}\,\dd x.
\]
Substitute $x=su$ ($\dd x = s\,\dd u$, limits $0$ to $1$):
\[
\int_0^s x^{\alpha_1-1}(s-x)^{\alpha_2-1}\dd x
 = s^{\alpha_1+\alpha_2-1}\int_0^1 u^{\alpha_1-1}(1-u)^{\alpha_2-1}\dd u
 = s^{\alpha_1+\alpha_2-1}\,B(\alpha_1,\alpha_2),
\]
where $B(\alpha_1,\alpha_2)=\Gamma(\alpha_1)\Gamma(\alpha_2)/\Gamma(\alpha_1+\alpha_2)$
is the Beta function -- the continuous counterpart of binomial coefficients
that normalizes polynomial powers on $[0,1]$. Substituting back, the Gamma functions telescope:
\[
f_S(s) = \frac{\beta^{\alpha_1+\alpha_2}}{\Gamma(\alpha_1+\alpha_2)}\,
         s^{\alpha_1+\alpha_2-1}e^{-\beta s},
\]
the $\Gam(\alpha_1+\alpha_2,\beta)$ density. The Beta integral is doing all
the work; this is the same combinatorics-of-the-continuum that made the
uniforms convolve into a tent.
\end{whybox}

\begin{corollary}[Exponentials sum to an Erlang]
\label{ch17:cor-erlang}
Since $\Expo(\beta)=\Gam(1,\beta)$, the sum of $n$ independent
$\Expo(\beta)$ variables is $\Gam(n,\beta)$ (the \emph{Erlang}
distribution). Concretely, the waiting time until the $n$th event in a
Poisson process of rate $\beta$ is $\Gam(n,\beta)$ -- you wait for $n$
independent exponential inter-arrival gaps and add them up.
\end{corollary}

\begin{applicationbox}
\textbf{Convolution in machine learning, two ways.} The same word shows up
in deep learning, and the connection is exact. A convolutional layer slides
a small kernel across an input and, at each position, computes
$\sum_k w_k\,x_{s-k}$ -- the discrete convolution of
Definition~\ref{ch17:def-discrete-conv} with the kernel $w$ playing the role
of one PMF. The probabilistic reading also pays off directly: when a system
adds independent random delays -- a request hops through several
microservices, each adding an $\Expo(\beta)$-like latency -- the end-to-end
latency is the convolution of the per-hop densities, an Erlang by
Corollary~\ref{ch17:cor-erlang}. Engineers who size timeouts from the
\emph{tail} of that convolved distribution (the $99.9$th percentile) rather
than the mean are using exactly this chapter's machinery to budget for the
rare slow request.
\end{applicationbox}

\begin{trapbox}
\textbf{The matching-parameter trap.} Each closure theorem comes with a
\emph{shared} parameter that you may not violate:
\begin{itemize}
  \item Binomials add only with the \textbf{same $p$}.
  \item Gammas add only with the \textbf{same rate $\beta$}.
  \item Normals add for \textbf{any} parameters (no matching needed) --
        that is part of what makes them special.
  \item Poissons add for any rates (the rates simply sum).
\end{itemize}
A common error is to ``add'' two binomials with different $p$, or two
gammas with different rates, and report a member of the same family. The
true sum exists, of course, but it is \emph{not} in the family and has no
tidy closed form. When the shared parameter does not match, you must fall
back to the raw convolution.
\end{trapbox}

\begin{labbox}
\textbf{Work it yourself.} Convince yourself the closures are mutually
consistent. Start from $X_i\sim\Bern(p)$ i.i.d.
\begin{enumerate}
  \item Each $X_i$ is $\Bin(1,p)$. Use Theorem~\ref{ch17:thm-bin-add}
        $n$ times to conclude $\sum_{i=1}^n X_i\sim\Bin(n,p)$. (This is
        the very \emph{definition} of the binomial -- nice that the
        convolution machinery reproduces it.)
  \item Each $\Expo(\beta)$ is $\Gam(1,\beta)$. Use
        Theorem~\ref{ch17:thm-gamma-add} to get the Erlang of
        Corollary~\ref{ch17:cor-erlang}.
  \item Each $\Pois(\lambda/n)$ is independent; their sum over $i=1,\dots,n$
        is $\Pois(\lambda)$ by Theorem~\ref{ch17:thm-pois-add}. So a
        Poisson of rate $\lambda$ is an aggregate of many tiny Poissons --
        the ``infinitely divisible'' property (meaning the distribution
        can be split into the sum of an arbitrary number $n$ of independent,
        identically distributed random pieces).
\end{enumerate}
\end{labbox}

\section{Expectation and Variance of Sums, Revisited}
\label{ch17:mean-var}

Convolution gives the \emph{entire} distribution of a sum, which is more
than we often need. Frequently we only want two numbers: the mean and the
variance. Those obey straightforward algebraic rules, and
the expectation rule needs \emph{no} independence at all.

\begin{theorem}[Linearity of expectation and additivity of variance]
\label{ch17:thm-mean-var}
For \emph{any} random variables $X_1,\dots,X_n$ with finite means,
\[
\E\!\left[\sum_{i=1}^n X_i\right] = \sum_{i=1}^n \E[X_i].
\]
For the variance, in general
\[
\Var\!\left(\sum_{i=1}^n X_i\right)
  = \sum_{i=1}^n \Var(X_i) + 2\!\!\sum_{1\le i<j\le n}\!\! \Cov(X_i,X_j),
\]
and if the $X_i$ are \emph{(pairwise) uncorrelated} -- in particular if
they are independent -- the covariance terms vanish:
\[
\Var\!\left(\sum_{i=1}^n X_i\right) = \sum_{i=1}^n \Var(X_i).
\]
\end{theorem}

\begin{happeningbox}
The contrast here is one of the most important lessons in the subject.
\emph{Means always add.} It does not matter whether your variables are
independent, correlated, tangled, adversarial -- expectation is linear,
full stop. \emph{Variances add only when there is no correlation.} If your
variables move together ($\Cov>0$), the sum swings more wildly than the
individual variances would suggest; if they move oppositely ($\Cov<0$),
they cancel and the sum is calmer. This single asymmetry -- mean is always
linear, spread is not -- underlies portfolio diversification, the design
of experiments, and the $\sqrt{n}$ that governs the law of large numbers.
\end{happeningbox}

\begin{focusbox}
\textbf{The $\sqrt{n}$ that runs through all of probability.}
Let $X_1,\dots,X_n$ be i.i.d.\ with mean $\mu$ and variance $\sigma^2$,
and form $S_n=\sum X_i$. Then
\[
\E[S_n]=n\mu, \qquad \Var(S_n)=n\sigma^2, \qquad \SD(S_n)=\sigma\sqrt{n}.
\]
The sum's \emph{typical size} grows like $n$, but its \emph{typical
wobble} grows only like $\sqrt n$. So the wobble \emph{relative to} the
total shrinks like $\sqrt n / n = 1/\sqrt n \to 0$. For the average
$\bar X_n=S_n/n$ this reads $\Var(\bar X_n)=\sigma^2/n$, with standard
deviation $\sigma/\sqrt n$. That vanishing $\sigma/\sqrt n$ is the engine
of the law of large numbers (the average concentrates on $\mu$) and, once
you rescale by exactly $\sqrt n$, the engine of the central limit theorem.
Remember the slogan: \textbf{sums grow like $n$, fluctuations like
$\sqrt n$.}
\end{focusbox}

\begin{applicationbox}
\textbf{Diversification in a portfolio.} Put your money equally in $n$
assets, each with return of variance $\sigma^2$. If the assets are
\emph{independent}, the portfolio return $\bar X_n$ has variance
$\sigma^2/n$ -- risk falls toward zero as you add names. This is the
mathematical content of ``don't put all your eggs in one basket.'' But the
\emph{covariance} terms are the catch: in a market crash assets fall
\emph{together} ($\Cov>0$), the cross terms in
Theorem~\ref{ch17:thm-mean-var} do not vanish, and diversification stalls.
The reason a diversified portfolio still carries irreducible
``systematic'' risk is precisely those non-zero covariances. The clean
$1/n$ decay is a fantasy that requires independence the real world does
not grant.
\end{applicationbox}

\begin{example}[Mean and variance without finding the full distribution]
\label{ch17:ex-meanvar}
You sum $n=100$ i.i.d.\ rolls of a fair six-sided die. A single die has
$\E[X]=3.5$ and $\Var(X)=\frac{35}{12}\approx 2.9167$. Without ever
computing the (genuinely messy) convolution of $100$ uniforms, we read off
\[
\E[S_{100}] = 100\cdot 3.5 = 350,\qquad
\Var(S_{100}) = 100\cdot\tfrac{35}{12} \approx 291.7,\qquad
\SD(S_{100}) \approx 17.1.
\]
So a total near $350\pm 17$ is typical; a total of $450$ would be about
$6$ standard deviations out -- effectively impossible. We extracted the
practical content of the distribution from two one-line rules, no
convolution required. (And by the central limit theorem the full
distribution is very nearly $\Normal(350,291.7)$.)
\end{example}

\section{A Paradox of Sums: The Cauchy That Refuses to Average}
\label{ch17:cauchy}

The examples above suggest that adding many independent variables and
dividing by $n$ causes the average to stabilize around a fixed number.
However, this stabilization requires the summands to have a well-defined
expectation. When that condition fails, the sample mean behaves very
differently.

\begin{paradoxbox}
\textbf{The standard Cauchy distribution} has density
$f(x)=\dfrac{1}{\pi(1+x^2)}$ on all of $\R$. It looks like a tame,
symmetric bell. But it has \emph{heavy tails} -- tails that decay so slowly
that extreme outliers far from the center retain substantial probability.
Because the density decays only like $1/x^2$ as $\abs{x}\to\infty$, the
integral $\int_{-\infty}^{\infty} |x|\,f(x)\,\dd x$ diverges to infinity.
The Cauchy has \textbf{no mean} (and therefore no variance).

Here is the surprise. If $X_1,\dots,X_n$ are i.i.d.\ standard Cauchy, their
\emph{average} $\bar X_n=\frac{1}{n}\sum X_i$ is again \emph{exactly}
standard Cauchy -- the \emph{same} distribution, for every $n$. Averaging
a million Cauchy samples is statistically no better than taking one. The
sample mean does not converge to anything; it just keeps wandering with
the same spread forever.
\end{paradoxbox}

\begin{figure}[H]
\centering
\begin{tikzpicture}
\begin{axis}[
  width=12cm, height=6cm,
  xlabel={$x$}, ylabel={density},
  xmin=-6, xmax=6, ymin=0, ymax=0.45,
  xtick={-6,-4,-2,0,2,4,6}, ytick={0,0.1,0.2,0.3,0.4},
  legend pos=north east, legend style={font=\scriptsize},
  axis lines=left, samples=120,
]
\addplot[densityred, thick, domain=-6:6] {1/sqrt(2*pi)*exp(-x^2/2)};
\addlegendentry{standard normal}
\addplot[distblue, thick, domain=-6:6] {1/(pi*(1+x^2))};
\addlegendentry{standard Cauchy}
\end{axis}
\end{tikzpicture}
\caption{The standard normal (red) and standard Cauchy (blue) both look like
symmetric bells near the center, but the Cauchy's tails decay only like
$1/x^2$ -- far slower than the normal's $e^{-x^2/2}$. Those fat tails carry
enough probability far out that $\int|x|\,f(x)\,\dd x$ diverges, so the Cauchy
has no mean. A single extreme draw can dominate any average, which is why
averaging Cauchy samples never settles down.}
\label{ch17:fig-cauchy}
\end{figure}

\begin{whybox}
How can the average not improve? Two facts collide. First, the Cauchy is
closed under addition like the normal: the convolution of two standard
Cauchys is a Cauchy of scale $2$, and in general the sum of $n$ standard
Cauchys is a Cauchy of scale $n$. (One slick proof: the characteristic
function of a standard Cauchy is $e^{-|t|}$; multiplying $n$ of them gives
$e^{-n|t|}$, the characteristic function of a Cauchy of scale $n$ --
exactly the convolution-becomes-multiplication trick of the next chapter.)
Dividing that sum by $n$ rescales by $1/n$, turning scale $n$ back into
scale $1$. So the average is standard Cauchy again.

Second, the law of large numbers \emph{requires a finite mean}. Its proof
leans on $\Var(\bar X_n)=\sigma^2/n\to 0$, and the Cauchy has neither a
finite variance nor even a finite mean to converge to. The $\sqrt n$
slogan silently assumed finite variance. When the tails are heavy enough
to destroy the mean, the whole averaging machine seizes up.
\end{whybox}

\begin{aestheticbox}
The Cauchy is the cautionary twin of the normal. Both are stable under
addition; both are fixed points of convolution. But the normal pulls
probability \emph{inward} and tames sums, while the Cauchy's fat tails (Figure~\ref{ch17:fig-cauchy}) let
a single freak sample dominate the whole average no matter how many other
samples you collect. The lesson for data science is blunt: before you
trust an average, ask whether the thing you are averaging even \emph{has}
a mean. Heavy-tailed data -- file sizes, city populations, financial
losses, network traffic -- can behave far more like Cauchy than like the
normal your instincts default to.
\end{aestheticbox}

\section{Setting Up the Limit Theorems}
\label{ch17:limit-setup}

The convolution machinery points straight toward the central limit
theorem. We have all the pieces; what remains is to see how they assemble.

\begin{intuitionbox}
Take i.i.d.\ $X_1,X_2,\dots$ with mean $\mu$ and finite variance
$\sigma^2$, and let $S_n=\sum_{i=1}^n X_i$. The PMF/PDF of $S_n$ is the
$n$-fold convolution of the common distribution with itself:
\[
f_{S_n} = f * f * \cdots * f \quad (n\ \text{copies}).
\]
We watched this in miniature: one uniform (flat) convolves into a triangle
($n{=}2$), then a bell ($n{=}3$), smoother and smoother. \emph{Convolving
a distribution with itself repeatedly drives it toward a single universal
shape.} Two clean facts pin that shape down:
\begin{itemize}
  \item \textbf{Center.} $\E[S_n]=n\mu$, so the bump marches off to the
        right (or left). Re-center by subtracting $n\mu$.
  \item \textbf{Spread.} $\Var(S_n)=n\sigma^2$, so the bump widens like
        $\sqrt n$. Re-scale by dividing by $\sigma\sqrt n$.
\end{itemize}
The exactly-tuned object
\[
Z_n=\frac{S_n-n\mu}{\sigma\sqrt n}
\]
has mean $0$ and variance $1$ for every $n$ -- centered and standardized
so it can neither run away nor blow up nor collapse. What shape does its
distribution converge to?
\end{intuitionbox}

\begin{figure}[H]
\centering
\begin{tikzpicture}
\begin{axis}[
  width=13cm, height=6cm,
  xlabel={value}, ylabel={density / scaled PMF},
  xmin=-0.1, xmax=4.1, ymin=0, ymax=1.15,
  xtick={0,1,2,3,4}, ytick={0,0.5,1},
  legend pos=north east, legend style={font=\scriptsize},
  axis lines=left, samples=100,
]
\addplot[gray, thick, domain=0:1] {1};
\addlegendentry{$n=1$ (flat)}
\addplot[gray, thick, domain=-0.1:0, forget plot] {0};
\addplot[gray, thick, domain=1:4.1, forget plot] {0};
\addplot[distblue, thick, domain=0:1] {x};
\addlegendentry{$n=2$ (triangle)}
\addplot[distblue, thick, domain=1:2, forget plot] {2-x};
\addplot[densityred, thick, domain=0:1] {x^2/2};
\addlegendentry{$n=3$ (bell-like)}
\addplot[densityred, thick, domain=1:2, forget plot] {(-2*x^2+6*x-3)/2};
\addplot[densityred, thick, domain=2:3, forget plot] {(3-x)^2/2};
\end{axis}
\end{tikzpicture}
\caption{Repeatedly convolving a $\Unif(0,1)$ density with itself. The flat
box ($n{=}1$) becomes a triangle ($n{=}2$) becomes a smooth, bell-shaped
curve ($n{=}3$). Each extra summand smooths and centers the shape. Recenter
and rescale, and the limit is the normal curve -- this picture \emph{is}
the central limit theorem in slow motion.}
\label{ch17:fig-clt}
\end{figure}

\begin{aestheticbox}
\textbf{The central limit theorem (preview).} For i.i.d.\ summands with
finite mean $\mu$ and variance $\sigma^2$,
\[
Z_n=\frac{S_n-n\mu}{\sigma\sqrt n}\ \dto\ \Normal(0,1)
\qquad\text{as } n\to\infty,
\]
\emph{no matter what distribution the $X_i$ came from} (uniform, Bernoulli,
exponential, your own homemade lump -- anything with finite variance).
The arrow $\dto$ denotes \emph{convergence in distribution}: as $n$ grows,
the cumulative distribution function of the standardized sum $Z_n$ approaches
the bell curve's CDF $\Phi(z)$ at every single point $z$.
This chapter built the gears: convolution is the operation that adds
independent variables; repeated convolution smooths any starting shape
(Figure~\ref{ch17:fig-clt}); linearity of expectation and additivity of
variance supply the exact centering $n\mu$ and scaling $\sigma\sqrt n$; and
the normal is the unique shape that convolution leaves unchanged, so it is
the only possible destination (Figure~\ref{ch17:fig-bell}). The Cauchy detour warned us why the
``finite variance'' clause cannot be dropped. A few chapters from now we
prove it.
\end{aestheticbox}

\begin{figure}[H]
\centering
\begin{tikzpicture}
\begin{axis}[
  width=12cm, height=5.8cm,
  xlabel={$z$}, ylabel={$\varphi(z)$},
  xmin=-4, xmax=4, ymin=0, ymax=0.46,
  xtick={-3,-2,-1,0,1,2,3}, ytick={0,0.1,0.2,0.3,0.4},
  axis lines=left, samples=120,
]
\addplot[name path=curve, densityred, thick, domain=-4:4]
  {1/sqrt(2*pi)*exp(-x^2/2)};
\path[name path=axis] (axis cs:-4,0) -- (axis cs:4,0);
\addplot[densityred!30] fill between[of=curve and axis, soft clip={domain=-4:-1.96}];
\addplot[densityred!30] fill between[of=curve and axis, soft clip={domain=1.96:4}];
\node at (axis cs:-2.7,0.05) [font=\scriptsize] {$2.5\%$};
\node at (axis cs:2.7,0.05) [font=\scriptsize] {$2.5\%$};
\draw[gray, dashed] (axis cs:-1.96,0) -- (axis cs:-1.96,0.4);
\draw[gray, dashed] (axis cs:1.96,0) -- (axis cs:1.96,0.4);
\node at (axis cs:0,0.43) [font=\scriptsize] {middle $95\%$};
\end{axis}
\end{tikzpicture}
\caption{The destination of the standardized sum $Z_n$: the standard normal
density $\varphi(z)=\frac{1}{\sqrt{2\pi}}e^{-z^2/2}$. The shaded tails beyond
$\pm 1.96$ each hold about $2.5\%$ of the probability, so roughly $95\%$ of
the mass lies within two standard deviations of the mean -- the rule of thumb
that powers confidence intervals once the central limit theorem licenses the
normal approximation.}
\label{ch17:fig-bell}
\end{figure}

\begin{connectionbox}
The leap from ``convolutions are hard integrals'' to ``the limit is
normal'' goes through \textbf{generating functions} and
\textbf{characteristic functions}, the subject of the very next chapter.
The key identity: if $\varphi_X(t)=\E[e^{itX}]$ is the characteristic
function, then independence turns the convolution of densities into the
\emph{product} of characteristic functions,
\[
\varphi_{X+Y}(t)=\varphi_X(t)\,\varphi_Y(t),
\]
and the $n$-fold convolution into the $n$th \emph{power}
$\varphi_{S_n}(t)=\varphi_X(t)^n$. Hard convolutions collapse into easy
algebra, and a Taylor expansion of $\varphi$ around $t=0$ -- whose
coefficients are exactly the mean and variance from this chapter -- makes
the Gaussian limit fall out almost by itself. Convolution is the
\emph{what}; the transform is the \emph{how}.
\end{connectionbox}

\section{Summary}
\label{ch17:summary}

\begin{formalbox}
\textbf{The distribution of a sum of independent variables.}
\begin{itemize}
  \item \textbf{Discrete:} $p_{X+Y}(s)=\sum_k p_X(k)\,p_Y(s-k)$.
  \item \textbf{Continuous:} $f_{X+Y}(s)=\int f_X(x)\,f_Y(s-x)\,\dd x$.
  \item Both are the \emph{convolution} $p_X*p_Y$ / $f_X*f_Y$: sum (or
        integrate) over every way the two pieces split the target total.
\end{itemize}
\textbf{Closure under addition} (independent summands):
\begin{itemize}
  \item $\Pois(\lambda)+\Pois(\mu)=\Pois(\lambda+\mu)$ (any rates).
  \item $\Bin(m,p)+\Bin(n,p)=\Bin(m+n,p)$ (same $p$).
  \item $\Normal(\mu_1,\sigma_1^2)+\Normal(\mu_2,\sigma_2^2)
        =\Normal(\mu_1+\mu_2,\sigma_1^2+\sigma_2^2)$ (any parameters).
  \item $\Gam(\alpha_1,\beta)+\Gam(\alpha_2,\beta)=\Gam(\alpha_1+\alpha_2,\beta)$
        (same rate); $n$ exponentials sum to an Erlang.
\end{itemize}
\textbf{Two numbers, two rules:} $\E[\sum X_i]=\sum\E[X_i]$ \emph{always};
$\Var(\sum X_i)=\sum\Var(X_i)$ \emph{when uncorrelated}. For i.i.d.\
summands, sums grow like $n$ and fluctuations like $\sqrt n$.
\textbf{Caveat:} heavy tails (Cauchy) can destroy the mean and break
averaging entirely.
\end{formalbox}

\subsection*{Looking Forward: Generating Functions}

Adding independent random variables corresponds to convolving their distributions.
While a few families have clean closed forms, repeated convolution integrals quickly
become unwieldy. Chapter~18 introduces \textbf{generating functions} (probability generating
functions, moment generating functions, and characteristic functions), which transform
convolution of distributions into simple multiplication of functions.

\section{Exercises}
\label{ch17:exercises}

\begin{exercise}[Mental warm-ups: quick sum intuition]
\label{ch17:ex-warmup}
Answer each of the following in your head in ten seconds:
\begin{enumerate}[label=(\alph*)]
  \item If $X \sim \Pois(3)$ and $Y \sim \Pois(4)$ are independent, what is
        the probability that their sum $X+Y$ equals zero?
  \item If $X \sim \Normal(10, 9)$ and $Y \sim \Normal(5, 16)$ are
        independent, what is the standard deviation of $X+Y$? What is the
        standard deviation of $X-Y$?
  \item You roll two fair six-sided dice. What is the most likely total, and
        how many of the $36$ equally likely outcomes produce that total?
\end{enumerate}
\end{exercise}

\begin{exercise}[Convolve by hand]
\label{ch17:ex-byhand}
Let $X\sim\Unif\set{1,2,3}$ and, independently, $Y$ take value $0$ with
probability $1/2$ and value $3$ with probability $1/2$. Use the discrete
convolution formula to write out the full PMF of $S=X+Y$. Verify the
probabilities sum to $1$, and compute $\E[S]$ two ways: directly from your
PMF, and via linearity of expectation.
\end{exercise}

\begin{exercise}[Poisson merging]
\label{ch17:ex-pois}
A help desk receives emails as a $\Pois(8)$ count per hour and phone calls
as an independent $\Pois(5)$ count per hour.
(a) What is the distribution of the total number of contacts per hour?
(b) Compute $\P(\text{total}=0)$ in one hour.
(c) Over an $8$-hour shift, what is the distribution and the expected total?
\end{exercise}

\begin{exercise}[Uniform sum density]
\label{ch17:ex-unif}
Let $X\sim\Unif(0,2)$ and $Y\sim\Unif(0,1)$ be independent. Carefully
carry out the convolution integral
$f_S(s)=\int f_X(x)f_Y(s-x)\dd x$, paying attention to the overlap of the
two boxes, and show that $f_S$ is a \emph{trapezoid}: it ramps up, stays
flat, then ramps down. Give the breakpoints and the flat height.
\end{exercise}

\begin{exercise}[Variance with correlation]
\label{ch17:ex-cov}
Let $X$ and $Y$ each have variance $4$, with $\Corr(X,Y)=\rho$.
(a) Express $\Var(X+Y)$ in terms of $\rho$.
(b) For which $\rho$ is $\Var(X+Y)$ largest? Smallest? Interpret both
extremes in words (perfectly aligned vs.\ perfectly opposing variables).
(c) For what value of $\rho$ does $X+Y$ have the same variance as a single
variable, namely $4$?
\end{exercise}

\begin{exercise}[Erlang waiting time]
\label{ch17:ex-erlang}
Buses arrive as a Poisson process at rate $\beta=3$ per hour, so
inter-arrival times are i.i.d.\ $\Expo(3)$.
(a) Name the distribution of the waiting time $T$ until the \emph{fourth}
bus, with its parameters.
(b) Compute $\E[T]$ and $\Var(T)$ using the additivity rules (do not
integrate).
(c) Explain in one sentence why $\SD(T)$ is smaller than four times the
$\SD$ of a single gap.
\end{exercise}

\begin{exercise}[The $\sqrt n$ law, numerically]
\label{ch17:ex-sqrtn}
A single fair coin scored $+1$ for heads and $-1$ for tails has mean $0$
and variance $1$. Let $S_n$ be the sum of $n$ such i.i.d.\ scores (a simple
random walk).
(a) Give $\E[S_n]$ and $\SD(S_n)$.
(b) For $n=10{,}000$, roughly how far from $0$ is $S_n$ typically? Is
$S_n=500$ surprising? Is $S_n=5000$?
(c) Now consider the \emph{fraction} of heads, $(S_n/n + 1)/2$. How does
its standard deviation behave as $n\to\infty$, and what theorem does that
foreshadow?
\end{exercise}

\begin{exercise}[Why same-$p$ matters]
\label{ch17:ex-difp}
Let $X\sim\Bin(1,0.2)$ and $Y\sim\Bin(1,0.8)$ be independent.
(a) Write the full PMF of $S=X+Y$ on $\set{0,1,2}$.
(b) If $S$ were $\Bin(2,p)$ for some $p$, the value $\P(S=2)$ would force a
particular $p$; find it, then check whether that same $p$ reproduces
$\P(S=0)$. Conclude that $S$ is \emph{not} binomial, and explain which
hypothesis of Theorem~\ref{ch17:thm-bin-add} fails.
\end{exercise}

\begin{exercise}[Normal closure and standardization]
\label{ch17:ex-normal}
Let $X\sim\Normal(3,4)$ and $Y\sim\Normal(-1,9)$ be independent.
(a) State the distribution of $X+Y$.
(b) State the distribution of $X-Y$ (be careful: the variance still adds).
(c) Find a constant $a\in[0,1]$ so that $aX+(1-a)Y$ has the smallest
possible variance, and report that minimum variance. (This is the seed of
inverse-variance weighting in statistics.)
\end{exercise}

\begin{exercise}[Proof: convolution is commutative]
\label{ch17:ex-commute}
Show directly from Definition~\ref{ch17:def-cont-conv} that
$f_X*f_Y=f_Y*f_X$, i.e.\ $\int f_X(x)f_Y(s-x)\dd x=\int f_Y(x)f_X(s-x)\dd x$,
using the substitution $u=s-x$. Then explain in one sentence why this is
obvious \emph{probabilistically}, with no integral at all (hint: what is
$X+Y$ versus $Y+X$?).
\end{exercise}

\begin{exercise}[Heavy tails break averaging]
\label{ch17:ex-cauchy}
Two independent standard Cauchy variables $X_1,X_2$ have
$\bar X_2=(X_1+X_2)/2$.
(a) Using the stated fact that $X_1+X_2$ is Cauchy of scale $2$, argue that
$\bar X_2$ is standard Cauchy again.
(b) Explain why this does not contradict the law of large numbers --
which hypothesis of that law does the Cauchy violate?
(c) In one sentence, what practical advice does this give a data analyst
facing heavy-tailed measurements?
\end{exercise}

\begin{exercise}[Convolving exponentials with unequal rates]
\label{ch17:ex-expo-diff}
Theorem~\ref{ch17:thm-gamma-add} showed that independent exponentials with
the \emph{same} rate $\beta$ convolve into an Erlang density $\Gam(2,\beta)$.
What happens when their rates differ? Let $X \sim \Expo(\lambda_1)$ and
$Y \sim \Expo(\lambda_2)$ be independent with $\lambda_1 \ne \lambda_2$.
\begin{enumerate}[label=(\alph*)]
  \item Write out the convolution integral
        $f_S(s) = \int_0^s \lambda_1 e^{-\lambda_1 x} \lambda_2 e^{-\lambda_2(s-x)}\,\dd x$
        for $s > 0$.
  \item Evaluate the integral and show that the density of the sum is
        \[
        f_S(s) = \frac{\lambda_1 \lambda_2}{\lambda_2 - \lambda_1}
                 \left( e^{-\lambda_1 s} - e^{-\lambda_2 s} \right)
        \qquad (s > 0).
        \]
  \item Use L'H\^opital's rule to evaluate the limit of $f_S(s)$ as
        $\lambda_2 \to \lambda_1$, and verify that you recover the Erlang
        density $\lambda_1^2 s e^{-\lambda_1 s}$ of
        Corollary~\ref{ch17:cor-erlang}.
\end{enumerate}
\end{exercise}
